\section{A Causal Test for Every Modality}
\label{sec:protocol}

Agreement with human labels cannot show that an attribution is faithful, because a model may answer from parts of the input that a person would not mark.
We therefore test attributions by intervention.

\paragraph{Replay with hidden units.}
We keep the model's own response $\mathbf{y}$ fixed and score it with teacher forcing, $\ell(\mathbf{y}\mid\mathbf{x})=\frac{1}{T}\sum_t\log P(y_t\mid\mathbf{x},y_{<t})$.
For a set of units $\mathcal{H}$, the \emph{drop} is $\Delta(\mathcal{H})=\ell(\mathbf{y}\mid\mathbf{x})-\ell(\mathbf{y}\mid\mathbf{x}_{\setminus\mathcal{H}})$, in nats per token, where $\mathbf{x}_{\setminus\mathcal{H}}$ is the input with the units in $\mathcal{H}$ hidden.

\paragraph{Hiding without a trace.}
A hidden unit is replaced \emph{at the input of the encoder} by a placeholder of the same length that carries no content: black frames for video, silence for audio, a black image, and blank tokens for a text passage.
The number of tokens, their positions, and the time stamps of all other units stay the same, and no content of the hidden unit reaches the model through the encoder.
In \Cref{appendix:masking} we repeat the tests with the mean-embedding replacement used in earlier work and reach the same conclusions on audio and video.

\paragraph{Metrics.}
All metrics compare a method's units with random units of the same count, drawn three times per question.
\emph{Lift} is the extra drop when the top-ranked units are hidden: $\Delta(\text{top})-\mathbb{E}\,\Delta(\text{random})$.
\emph{Keep gain} is the reverse test: we hide everything except the selected units, and report how much less the response loses than when the same number of random units is kept.
At sentence level we hide the top units of sentence $C_k$ and measure the drop of that sentence alone. \emph{Sentence lift} compares it with random units, and \emph{sentence specificity} compares it with the mean drop of the other sentences of the same response, so it is positive only if the units matter for this sentence in particular.
The budget is a share of the duration for audio and video (10, 20, 30, and 50\%) and a number of quotes for documents (1, 2, or 3 of 15). The main tables use the middle budget, 20\% or 2 quotes; sentence-level tests use 20\% or 1 quote.
We report means over questions with 95\% bootstrap intervals, and compare methods by paired differences over the same questions.

\section{Experiments}
\label{sec:exp}

\begin{table}[t]
\centering
\tablesize
\setlength{\tabcolsep}{2.2pt}
\begin{tabular}{lccc}
\toprule
 & \textbf{Video} & \textbf{Audio} & \textbf{Image--text}\\
\midrule
Dataset & NExT-GQA & MMAU & MMDocRAG \\
Task & multiple choice & multiple choice & long answer \\
Questions (test) & 300 (246) & 249 (196) & 400 (400) \\
Source unit & 1 s of frames & 1 s of audio & passage / image \\
Units per question & 46 & 22 & 15 \\
Input length & 10--90 s & 1.5--35 s & 15 quotes \\
Human evidence & time spans & -- & gold quotes \\
\midrule
\multicolumn{4}{l}{\textit{Answer accuracy (\%)}}\\
Qwen2.5-Omni-7B & 72.3 & 68.3 & 23.8$^\ast$ \\
MiniCPM-o 4.5 & 77.0 & 76.3 & 50.5$^\ast$ \\
Qwen3-Omni-30B & 78.7 & 73.1 & 54.0$^\ast$ \\
\bottomrule
\end{tabular}
\caption{Data. Test questions exclude the split used to check the label threshold. $^\ast$Long answers graded as fully correct by an LLM judge against the reference answer (\Cref{appendix:judge}); counting partly correct answers gives 66.8, 85.8, and 86.8.}
\label{tab:stats}
\skiptableend
\end{table}


\begin{table*}[t]
\centering
\small
\setlength{\tabcolsep}{4.2pt}
% generated by analysis/paper_models.py; do not edit by hand
\begin{tabular}{lrcccrcccrccc}
\toprule
 & \multicolumn{4}{c}{\textbf{Video}} & \multicolumn{4}{c}{\textbf{Audio}} & \multicolumn{4}{c}{\textbf{Image--text}} \\
\cmidrule(lr){2-5}\cmidrule(lr){6-9}\cmidrule(lr){10-13}
\textbf{Method} & Fwd & Lift & S-lift & S-spec & Fwd & Lift & S-lift & S-spec & Fwd & Lift & S-lift & S-spec \\
\midrule
Embedding similarity & 0 & 0.008 & 0.009 & 0.005 & 0 & 0.034 & 0.031 & 0.005 & 0 & 0.216 & 0.216 & 0.156 \\
Self-citation & 1 & 0.013 & 0.012 & 0.009 & 1 & 0.040 & 0.036 & 0.024 & 1 & 0.184 & 0.349 & \underline{0.317} \\
Leave-one-out$^\ast$ & 41 & 0.016 & 0.024 & 0.022 & 17 & 0.100 & \underline{0.133} & \underline{0.105} & 16 & 0.368 & 0.404 & 0.286 \\
ContextCite$^\ast$ & 33 & 0.016 & 0.027 & \underline{0.023} & 28 & 0.099 & \textbf{0.139} & \textbf{0.119} & 33 & \underline{0.389} & 0.397 & 0.285 \\
\midrule
\multicolumn{13}{>{\columncolor{gray!15}}l}{\textit{OmniTrace on Qwen2.5-Omni-7B}}\\
\textbf{OmniTrace} (ours) & 0 & 0.016 & 0.015 & 0.007 & 0 & 0.077 & 0.110 & 0.084 & 0 & 0.351 & 0.376 & 0.245 \\
\textbf{OmniTrace-trained} (ours)$^\P$ & 0 & 0.014 & 0.016 & 0.010 & 0 & 0.081 & 0.116 & 0.086 & 0 & 0.374 & 0.385 & 0.250 \\
\midrule
\multicolumn{13}{>{\columncolor{gray!15}}l}{\textit{OmniTrace on MiniCPM-o 4.5}}\\
\textbf{OmniTrace} (ours) & 0 & 0.024 & 0.027 & 0.012 & 0 & \underline{0.136} & 0.102 & 0.047 & 0 & 0.105 & 0.159 & 0.098 \\
\textbf{OmniTrace-trained} (ours)$^\P$ & 0 & \textbf{0.067} & \textbf{0.065} & \textbf{0.024} & 0 & \textbf{0.169} & 0.120 & 0.054 & 0 & \textbf{0.432} & \textbf{0.706} & \textbf{0.507} \\
\midrule
\multicolumn{13}{>{\columncolor{gray!15}}l}{\textit{OmniTrace on Qwen3-Omni-30B-A3B}}\\
\textbf{OmniTrace} (ours) & 0 & \underline{0.033} & \underline{0.032} & 0.011 & 0 & 0.112 & 0.095 & 0.065 & 0 & 0.316 & \underline{0.448} & 0.309 \\
\bottomrule
\end{tabular}

\caption{\textbf{Causal faithfulness and cost on Qwen2.5-Omni-7B.}
Fwd: extra forward passes per question beyond the one generation.
Lift: extra drop in the response's log-probability (nats per token) when the method's top units are hidden, over hiding the same number of random units (20\% of the duration; 2 quotes).
S-lift and S-spec: sentence lift and sentence specificity.
\mdn / \mup: reliably lower / higher than OmniTrace on the same questions (paired 95\% bootstrap interval).
$^\dagger$Self-citation is scored on the responses of a separate generation and only where the model cites (62, 72, and 328 responses); it hides the cited units, and the row below it is OmniTrace on the same responses with the same number of units.
$^\ddagger$One extra generation.}
\label{tab:baselines}
\skiptableend
\end{table*}

\begin{table*}[t]
\vspace{4pt}
\centering
\small
\setlength{\tabcolsep}{4.2pt}
\resizebox{\textwidth}{!}{% generated by analysis/paper_models.py; do not edit by hand
\begin{tabular}{lcccccccccccc}
\toprule
 & \multicolumn{4}{c}{\textbf{Video}} & \multicolumn{4}{c}{\textbf{Audio}} & \multicolumn{4}{c}{\textbf{Image--text}} \\
\cmidrule(lr){2-5}\cmidrule(lr){6-9}\cmidrule(lr){10-13}
 & Lift & Keep & S-lift & S-spec & Lift & Keep & S-lift & S-spec & Lift & Keep & S-lift & S-spec \\
\midrule
\multicolumn{13}{>{\columncolor{gray!15}}l}{\textit{Qwen2.5-Omni-7B} \;($n$ = 246, 196, 400)}\\
\textbf{OmniTrace} & 0.016 & 0.010 & 0.015 & 0.007 & 0.069 & 0.066 & 0.103 & 0.084 & 0.348 & 0.412 & 0.375 & 0.245 \\
\;argmax tracing + curation & 0.009\mdn & 0.008 & 0.010\mdn & 0.004\mdn & 0.067 & 0.064 & 0.096\mdn & 0.081 & 0.346 & 0.420 & -- & -- \\
\;last layer only & 0.014 & 0.010 & -- & -- & 0.048\mdn & 0.046\mdn & -- & -- & 0.306\mdn & 0.380\mdn & -- & -- \\
\;per-token mean & -- & -- & -- & -- & -- & -- & -- & -- & 0.224\mdn & 0.255\mdn & 0.286\mdn & 0.214\mdn \\
Human-evidence units & 0.009\mdn & -- & -- & -- & -- & -- & -- & -- & 0.211\mdn & -- & -- & -- \\
\midrule
\multicolumn{13}{>{\columncolor{gray!15}}l}{\textit{MiniCPM-o 4.5} \;($n$ = 233, 187, 400)}\\
\textbf{OmniTrace} & 0.027 & 0.026 & 0.029 & 0.013 & 0.143 & 0.106 & 0.120 & 0.055 & 0.105 & 0.074 & 0.158 & 0.098 \\
\;argmax tracing + curation & 0.027 & $-$0.051\mdn & 0.025 & $-$0.001\mdn & 0.020\mdn & 0.001\mdn & 0.021\mdn & 0.006\mdn & 0.079\mdn & 0.058\mdn & -- & -- \\
\;last layer only & 0.025 & 0.005\mdn & -- & -- & 0.091\mdn & 0.055\mdn & -- & -- & 0.092 & 0.056 & -- & -- \\
\;per-token mean & -- & -- & -- & -- & -- & -- & -- & -- & 0.083 & 0.057 & 0.188\mup & 0.165\mup \\
Human-evidence units & 0.043\mup & -- & -- & -- & -- & -- & -- & -- & 0.334\mup & -- & -- & -- \\
\midrule
\multicolumn{13}{>{\columncolor{gray!15}}l}{\textit{Qwen3-Omni-30B-A3B} \;($n$ = 196, 400)}\\
\textbf{OmniTrace} & -- & -- & -- & -- & 0.112 & 0.097 & 0.095 & 0.065 & 0.316 & 0.304 & 0.448 & 0.309 \\
\;argmax tracing + curation & -- & -- & -- & -- & 0.112 & 0.105\mup & 0.087\mdn & 0.068 & 0.323 & 0.322\mup & -- & -- \\
\;last layer only & -- & -- & -- & -- & 0.045\mdn & 0.030\mdn & -- & -- & 0.125\mdn & 0.124\mdn & -- & -- \\
\;per-token mean & -- & -- & -- & -- & -- & -- & -- & -- & 0.240\mdn & 0.233\mdn & 0.381\mdn & 0.297\mdn \\
Human-evidence units & -- & -- & -- & -- & -- & -- & -- & -- & 0.267\mdn & -- & -- & -- \\
\bottomrule
\end{tabular}
}
\caption{\textbf{OmniTrace on two models, its variants, and human-marked evidence.}
Lift and Keep: lift and keep gain at response level (20\% of the duration; 2 quotes); S-lift and S-spec: sentence lift and specificity. All in nats per token.
\mdn / \mup: reliably lower / higher than OmniTrace on the same questions or sentences.
\emph{Argmax tracing + curation} traces each token to its highest-attention position and filters the votes; \emph{last layer only} and \emph{per-token mean} change one choice in \Cref{eq:mass}.
\emph{Human-evidence units} hides the units annotators marked (not available for audio); their number differs from the budget, so their lift is against random sets of their own size.
Nothing was tuned for MiniCPM-o~4.5.}
\label{tab:models}

\end{table*}


\subsection{Setup}

\paragraph{Models and data.}
We use Qwen2.5-Omni-7B~\cite{xu2025qwen2} and MiniCPM-o~4.5~\cite{yu2025minicpm} with greedy decoding and at most 512 new tokens.
\Cref{tab:stats} lists the data.
For video we use NExT-GQA~\cite{xiao2024nextgqa}, multiple-choice questions with human-marked evidence intervals; the model sees the frames only.
For audio we use the questions of MMAU~\cite{sakshi2025mmau} whose answer depends on a part of the clip, not on the clip as a whole.
For image--text input we use MMDocRAG~\cite{dong2025mmdocrag}: each question comes with 15 quotes from a long document, about ten text passages and five images (figures, tables, charts), of which about 19\% are marked as gold evidence.
The models are asked to reason step by step, so responses have 6 to 10 sentences on average.
A fifth of the video and audio questions was used to check the label threshold and is excluded from all results.

\paragraph{Baselines.}
Every baseline ranks the same units and goes through the same causal test.
\emph{Embedding similarity} scores a unit by the cosine similarity between a sentence and the unit, using SigLIP for frames and images, CLAP for audio, and a text encoder for passages.
\emph{Self-citation} repeats the generation with a prompt that asks the model to cite time spans or quote ids after each sentence.
\emph{Leave-one-out} hides each unit in turn and scores it by the drop, which costs one forward pass per unit.
\emph{ContextCite}~\cite{cohen-wang2024contextcite} fits a sparse linear model on 32 random ablations.
The two perturbation baselines use our placeholders, so they run on audio and video as they do on text; they are fitted with the same hiding operation and the same target that the test measures, which favors them.
% AT2: add one sentence here once its rows are in Table 2 (learned head weights, cross-fitted on the ContextCite ablations).
Baselines are run on Qwen2.5-Omni-7B. Details are in \Cref{appendix:baselines}.

\subsection{Comparison with Baselines}
\label{sec:baselines}


\Cref{tab:baselines} compares OmniTrace with the baselines.

\paragraph{Same cost: embedding similarity and self-citation.}
OmniTrace has about twice the lift of embedding similarity in all three modalities (0.016 against 0.008 on video, 0.077 against 0.034 on audio, 0.351 against 0.216 on image--text), and the gap is larger at sentence level.
Embedding similarity finds units that look related to the sentence, which is not the same as the units the model used: its sentence specificity on audio is close to zero.
Self-citation is discussed in \Cref{sec:stated}; where the model cites at all, the cited units are less effective than OmniTrace's on audio and image--text and equally effective on video.

\paragraph{Higher cost: perturbation.}
Leave-one-out and ContextCite need 16 to 41 extra forward passes per question.
At response level OmniTrace reaches 77\% and 78\% of their lift on audio, 95\% and 90\% on image--text, and is level with both on video.
Over all budgets (\Cref{appendix:budgets}), the share is 67--90\% on audio and 88--100\% on image--text; on video OmniTrace is behind at the 10\% budget and level from 20\% on.
At sentence level the picture is the same on audio and image--text (79--95\%), and weaker on video, where OmniTrace reaches about 60\% of the perturbation baselines' sentence lift.
Since these baselines measure the tested quantity directly, we read them as a practical upper bound: OmniTrace recovers most of it from a single pass.

\subsection{A Second Model and Simpler Variants}
\label{sec:variants}


\Cref{tab:models} reports OmniTrace on both models next to variants that change one design choice.

\paragraph{The selected units matter, on both models.}
The lift of OmniTrace is above zero for both models, in every modality and at every budget (\Cref{appendix:budgets}). Keep gains are positive throughout, and reliably so except at some video budgets. All sentence lifts and specificities are above zero: hiding the units attributed to a sentence hurts that sentence more than it hurts the others.
The effect grows from video to audio to image--text. On video it is small in absolute terms: hiding the top 20\% of the seconds costs Qwen2.5-Omni 0.020 nats per token, about 4\% of the response's log-probability, because neighboring frames carry much of the same content.
The settings were chosen on Qwen2.5-Omni video data and applied to MiniCPM-o~4.5 unchanged; lifts there are larger on video and audio and smaller on image--text.

\paragraph{Hard tracing does not transfer.}
On Qwen2.5-Omni, argmax tracing with curation is level with OmniTrace on audio and image--text and behind on video.
On MiniCPM-o~4.5 it loses most of its effect on audio (lift 0.02 against 0.14, sentence lift 0.02 against 0.12), is behind on image--text, and on video its top units are worse to keep than random ones.
Its extra steps fit one model; the plain sum of attention mass does not need them.
Using the last layer alone is reliably worse on audio for both models, and the per-token mean is worse at response level on Qwen2.5-Omni (\Cref{sec:size}).

\paragraph{Labels.}
The labels of \Cref{eq:label} select 10--21\% of the units of a response.
Hiding the labelled units gives a positive lift in all six settings, and labelled units are specific to their sentence in all six (\Cref{appendix:labels}).
On image--text input the labels agree with the gold quotes at precision 0.48 and recall 0.60, where random selection has precision 0.19.

\paragraph{Human evidence is not always what the model uses.}
On Qwen2.5-Omni, the units OmniTrace selects have a larger lift than the units annotators marked as evidence (0.348 against 0.211 on image--text).
On MiniCPM-o~4.5 the order is reversed: hiding the human-marked quotes costs 0.334, three times OmniTrace's lift.
On that model OmniTrace finds units that matter, but not the ones that matter most; we return to this in the limitations.
